\chapter{Graph Theory — Union Find}
%%% generated by the blueprint pipeline from TCSlib.GraphTheory.Kruskal.UnionFind
\section{Overview}

This module establishes the correctness of the union--find data structure used by
Kruskal's algorithm: the relation $\mathtt{SamePartition}$ is an equivalence, merging
along a list of edges is invariant under permutation of that list, and two vertices end
up in the same component after merging all edges exactly when they are connected by a
path in the corresponding edge set.

%%%%%%%%%%%%%%%%%%%%%%%%%%%%%%%%%%%%%%%%%%%%%%%%%%%%%%%%%%%%%%%%%%%%%%%%%%%%%%%
\section{Declarations}

\begin{lemma}[Reflexivity of the same-partition relation]
\lean{Kruskal.UF.SamePartition.rfl}
\label{Kruskal.UF.SamePartition.rfl}
\leanok
\uses{Kruskal.UF, Kruskal.UF.SamePartition}
\difficulty{1}
Fix $n \in \bbn$, and recall that a union-find state on $n$ nodes is a function
$\mathrm{Fin}\,n \to \bbn$ assigning to each node a representative label. Every such
state $\mathit{uf}$ induces the same partition as itself: for all nodes $i, j :
\mathrm{Fin}\,n$, one has $\mathit{uf}(i) = \mathit{uf}(j)$ if and only if
$\mathit{uf}(i) = \mathit{uf}(j)$.
\end{lemma}

\begin{lemma}[Symmetry of the same-partition relation]
\lean{Kruskal.UF.SamePartition.symm}
\label{Kruskal.UF.SamePartition.symm}
\leanok
\uses{Kruskal.UF, Kruskal.UF.SamePartition}
\difficulty{1}
Fix $n \in \bbn$, and regard a union-find state on $n$ nodes as a function
$\mathrm{Fin}\,n \to \bbn$ assigning each node a representative. If two such states
$\mathit{uf}_1$ and $\mathit{uf}_2$ induce the same partition — that is, for all nodes
$i, j$ one has $\mathit{uf}_1(i) = \mathit{uf}_1(j)$ if and only if $\mathit{uf}_2(i) =
\mathit{uf}_2(j)$ — then $\mathit{uf}_2$ and $\mathit{uf}_1$ likewise induce the same
partition.
\end{lemma}

\begin{lemma}[Transitivity of the same-partition relation]
\lean{Kruskal.UF.SamePartition.trans}
\label{Kruskal.UF.SamePartition.trans}
\leanok
\uses{Kruskal.UF, Kruskal.UF.SamePartition}
\difficulty{1}
Let $n$ be a natural number, and let $\mathit{uf}_1, \mathit{uf}_2, \mathit{uf}_3 :
\mathrm{Fin}\,n \to \bbn$ be three union-find states on $n$ nodes. Suppose that
$\mathit{uf}_1$ and $\mathit{uf}_2$ induce the same partition of the nodes, and that
$\mathit{uf}_2$ and $\mathit{uf}_3$ induce the same partition. Then $\mathit{uf}_1$ and
$\mathit{uf}_3$ induce the same partition: for all nodes $i, j$, one has
$\mathit{uf}_1(i) = \mathit{uf}_1(j)$ if and only if $\mathit{uf}_3(i) =
\mathit{uf}_3(j)$.
\end{lemma}

\begin{lemma}[Merging within a single component leaves the state unchanged]
\lean{Kruskal.merge_noop}
\label{Kruskal.merge_noop}
\leanok
\uses{Kruskal.UF, Kruskal.UF.merge}
\difficulty{2}
Let $\mathit{uf}\colon \mathrm{Fin}\,n \to \bbn$ be a union-find state on $n$ nodes, and
let $i,j$ be two nodes whose representatives agree, $\mathit{uf}(i) = \mathit{uf}(j)$.
Then merging the components of $i$ and $j$ has no effect: the merged state equals
$\mathit{uf}$.
\end{lemma}

\begin{lemma}[Merging edges preserves shared components]
\lean{Kruskal.mergeAll_preserves}
\label{Kruskal.mergeAll_preserves}
\leanok
\uses{Kruskal.UF, Kruskal.UF.merge, Kruskal.UF.mergeAll, Kruskal.WEdge}
\difficulty{3}
Let $\mathit{uf} \colon \mathrm{Fin}\,n \to \bbn$ be a union-find state assigning each
of $n$ nodes a representative, and let $\mathit{es}$ be a list of weighted edges over
these $n$ nodes. Suppose two nodes $a, b \in \mathrm{Fin}\,n$ satisfy $\mathit{uf}(a) =
\mathit{uf}(b)$. Folding every edge of $\mathit{es}$ into $\mathit{uf}$ in sequence
yields a union-find state $\mathit{uf}'$ for which $\mathit{uf}'(a) = \mathit{uf}'(b)$;
that is, sequential merging never separates two nodes that already carry a common
representative.
\end{lemma}

\begin{lemma}[Merging preserves the reachability invariant]
\lean{Kruskal.mergeAll_iff_reach}
\label{Kruskal.mergeAll_iff_reach}
\leanok
\uses{Kruskal.Reach, Kruskal.UF, Kruskal.UF.merge, Kruskal.UF.mergeAll, Kruskal.WEdge, Kruskal.reach_cons_iff, Kruskal.reach_mono, Kruskal.reach_symm}
\difficulty{6}
Fix $n$ and let $uf : \mathrm{Fin}\,n \to \bbn$ be a union-find state assigning each of
the $n$ nodes a representative. Let $\mathit{base}$ and $\mathit{edges}$ be lists of
weighted edges on these $n$ vertices, and suppose $uf$ already encodes reachability
along $\mathit{base}$, in the sense that for all nodes $a, b$ one has $uf(a) = uf(b)$ if
and only if $a$ and $b$ are connected by a path through the edges of $\mathit{base}$.
Then, writing $uf'$ for the state obtained by merging the endpoints of every edge of
$\mathit{edges}$ into $uf$ in order, for all nodes $a, b$ one has $uf'(a) = uf'(b)$ if
and only if $a$ and $b$ are connected by a path through the edges of the concatenated
list $\mathit{base} \mathbin{+\!\!+} \mathit{edges}$.
\end{lemma}

\begin{lemma}[Correctness of union-find merging from the initial state]
\lean{Kruskal.mergeAll_init_iff}
\label{Kruskal.mergeAll_init_iff}
\leanok
\uses{Kruskal.Reach, Kruskal.SymAdj, Kruskal.UF.init, Kruskal.UF.mergeAll, Kruskal.WEdge, Kruskal.mergeAll_iff_reach}
\difficulty{4}
Fix a natural number $n$, and let $\mathit{edges}$ be a list of weighted edges on the
vertex set $\mathrm{Fin}\,n$. Consider the union-find state obtained by starting from
the initial state on $n$ nodes, in which every node is its own representative, and
processing $\mathit{edges}$ in order, so that at each weighted edge the components of
the two endpoints of that edge are merged. Then, for any two vertices $a, b \in
\mathrm{Fin}\,n$, the vertices $a$ and $b$ carry the same representative in the
resulting state if and only if $a$ and $b$ are connected by a path through the edges of
$\mathit{edges}$ — equivalently, if and only if the reflexive--transitive closure of
symmetric adjacency along $\mathit{edges}$ relates $a$ to $b$.
\end{lemma}

\begin{lemma}[Merging edges preserves equal partitions]
\lean{Kruskal.mergeAll_congr}
\label{Kruskal.mergeAll_congr}
\leanok
\uses{Kruskal.UF, Kruskal.UF.SamePartition, Kruskal.UF.merge, Kruskal.UF.mergeAll, Kruskal.WEdge}
\difficulty{3}
Let $\mathit{uf}_1, \mathit{uf}_2 : \mathrm{Fin}\,n \to \bbn$ be two union-find states
over $n$ nodes that induce the same partition, meaning that for all nodes $i, j$ one has
$\mathit{uf}_1(i) = \mathit{uf}_1(j)$ if and only if $\mathit{uf}_2(i) =
\mathit{uf}_2(j)$. Then for any list of weighted edges, folding that list into each
state by successively merging the endpoints of every edge yields two states that again
induce the same partition.
\end{lemma}

\begin{lemma}[Merge order is irrelevant to the induced partition]
\lean{Kruskal.merge_swap_partition}
\label{Kruskal.merge_swap_partition}
\leanok
\uses{Kruskal.UF, Kruskal.UF.SamePartition, Kruskal.UF.merge, Kruskal.WEdge}
\difficulty{3}
Let $n$ be a natural number, let $\mathit{uf}\colon \mathrm{Fin}\,n \to \bbn$ be a
union-find state, and let $a$ and $b$ be weighted edges over $n$ vertices, with
endpoints $a_u, a_v$ and $b_u, b_v$ respectively. Form one state by merging the
components of $a_u$ and $a_v$ and then merging the components of $b_u$ and $b_v$, and a
second state by performing these two merges in the opposite order. Then the two
resulting states induce the same partition of the nodes: for all $i, j :
\mathrm{Fin}\,n$, the first state assigns $i$ and $j$ the same representative if and
only if the second state does.
\end{lemma}

\begin{lemma}[Permutation invariance of edge merging]
\lean{Kruskal.mergeAll_perm}
\label{Kruskal.mergeAll_perm}
\leanok
\uses{Kruskal.UF, Kruskal.UF.SamePartition, Kruskal.UF.SamePartition.trans, Kruskal.UF.mergeAll, Kruskal.WEdge, Kruskal.erase_append_perm, Kruskal.mergeAll_congr, Kruskal.merge_swap_partition}
\difficulty{4}
Fix $n$ and an initial union-find state $\mathit{uf} : \mathrm{Fin}\,n \to \bbn$. If two
lists of weighted edges $l_1$ and $l_2$ over $n$ vertices are permutations of one
another, then folding each list into $\mathit{uf}$ by merging the endpoints of every
edge in turn yields two union-find states that induce the same partition of the nodes;
that is, for all $i, j : \mathrm{Fin}\,n$, the two states assign $i$ and $j$ equal
representatives under one exactly when they do under the other.
\end{lemma}

\begin{lemma}[Appending a final edge to a merge fold]
\lean{Kruskal.mergeAll_append_single}
\label{Kruskal.mergeAll_append_single}
\leanok
\uses{Kruskal.UF, Kruskal.UF.merge, Kruskal.UF.mergeAll, Kruskal.WEdge}
\difficulty{2}
Fix $n \in \bbn$ and a union-find state $\mathit{uf} \colon \mathrm{Fin}\,n \to \bbn$ on
$n$ nodes. For any list $l$ of weighted edges over $n$ vertices and any single weighted
edge $e$ over $n$ vertices with endpoints $u$ and $v$, merging all edges of the list $l
\mathbin{+\!\!+} [e]$ into $\mathit{uf}$ in order produces the same union-find state as
first merging all edges of $l$ into $\mathit{uf}$ in order and then merging the
components of $u$ and $v$ in the state so obtained.
\end{lemma}

\begin{lemma}[Peeling the leading edge off a merge fold]
\lean{Kruskal.mergeAll_cons_eq}
\label{Kruskal.mergeAll_cons_eq}
\leanok
\uses{Kruskal.UF, Kruskal.UF.SamePartition, Kruskal.UF.SamePartition.rfl, Kruskal.UF.merge, Kruskal.UF.mergeAll, Kruskal.WEdge}
\difficulty{1}
Let $\mathit{uf} : \mathrm{Fin}\,n \to \bbn$ be a union-find state on $n$ nodes, let $e$
be a weighted edge with endpoints $u$ and $v$ over these nodes, and let $S$ be a list of
weighted edges over the same nodes. Folding the list $e :: S$ into $\mathit{uf}$ and
folding the list $S$ into the union-find state that merges the components of $u$ and $v$
in $\mathit{uf}$ induce the same partition of the nodes: for all nodes $i$ and $j$,
folding $e :: S$ into $\mathit{uf}$ assigns $i$ and $j$ a common representative if and
only if folding $S$ into the merge of $u$ and $v$ assigns $i$ and $j$ a common
representative.
\end{lemma}

\begin{lemma}[Moving a member to the end is a permutation]
\lean{Kruskal.erase_append_perm}
\label{Kruskal.erase_append_perm}
\leanok
\uses{Kruskal.WEdge}
\difficulty{1}
Let $l$ be a finite list of weighted edges over $n$ vertices, and let $e$ be a weighted
edge that occurs in $l$. Then $l$ is a permutation of the list obtained by deleting the
first occurrence of $e$ from $l$ and appending $e$ at the end.
\end{lemma}
